% Preamble
\documentclass[../main.tex]{subfiles}

% Document
\begin{document}
\chapter{Conclusions}\label{ch:conclusions}

This document has presented a study of the multi-armed bandit problem and its extension 
to contextual multi-armed bandits, providing a detailed explanation of the main algorithms 
traditionally used to solve them.  

These types of problems, as shown in this work, can be transferred from the theoretical 
domain to the practical domain in multiple areas, and are therefore highly useful.  

In the case of using these algorithms for the implementation of recommendation systems, 
this represents a clear improvement over making recommendations in a completely random 
manner, which explains the rise of such systems in e-commerce platforms, social networks, 
and others.  

Regarding the algorithms implemented in the designed recommendation system, they either 
ignore context entirely or model it as a linear function, which in many practical cases 
does not accurately reflect reality.  

Nowadays, the state of the art in recommendation systems is increasingly shifting towards 
the use of more complex deep learning models. However, the use of these classical models 
will always have a place in certain applications, especially those where the available 
context is very limited or nonexistent, situations in which deep learning struggles to 
achieve good performance. For this reason, the study of these techniques remains of vital 
importance today.  

As for deep learning techniques that can be used to solve the multi-armed bandit problem, 
their study and performance comparison with the techniques analyzed in this document is 
left for future work.  
\end{document}